前面从加噪 SDE 推出概率流 ODE，是先规定随机过程，再导出确定性速度。
现在换一个起点：直接规定数据与噪声之间的插值，再学习能搬运这组边缘分布的速度。
这里保留 $t=0$ 为数据、$t=1$ 为噪声的约定。
取 $X_0\sim p_{\rm data}$ 与 $Z\sim\mathcal N(0,I)$ 独立，令
\[
X_t=(1-t)X_0+tZ,\qquad \dot X_t=Z-X_0.
\]
训练时两端点都已知，因此可以直接回归
\[
\mathcal L_{\rm CFM}(\theta)=
 \mathbb E_{t,X_0,Z}\|v_\theta(X_t,t)-(Z-X_0)\|^2.
\]
但生成时并不知道当前点对应哪一对 $X_0,Z$，所以模型实际学到的是条件平均
\[
v^*(x,t)=\mathbb E[Z-X_0\mid X_t=x].
\]
为什么这个条件平均仍能生成所需分布？对任意光滑测试函数 $\varphi$，
先对插值求导，再使用条件期望：
\begin{align}
\frac{\mathrm d}{\mathrm dt}\mathbb E[\varphi(X_t)]
 &=\mathbb E[\nabla\varphi(X_t)^\top(Z-X_0)]\\
 &=\int p_t(x)\nabla\varphi(x)^\top v^*(x,t)\,\mathrm dx.
\end{align}
分部积分后得到 $\partial_tp_t=-\nabla\cdot(p_tv^*)$。
因此在适当存在唯一性与边界条件下，这个速度场的 ODE 具有插值指定的逐时边缘分布。
训练时无需先数值求解 ODE 来获得监督；采样时则从噪声端反向积分。
这就是条件 Flow Matching 将困难的边缘速度回归改成已知条件标签回归的原因。
\cite{fm,rf,flow-guide}

这里的“直线”首先是每对训练端点之间的条件插值。
经过同一个 $x$ 的端点对可能很多，条件平均后的速度会随空间与时间改变，
所以由它积分出来的轨迹未必都是直线，更不保证一个 Euler 步就到达真实数据。
例如上面模型在 $t=1$ 处有 $X_1=Z$，且
$v^*(z,1)=z-\mathbb E[X_0]$。若直接从1走到0做一个 Euler 大步，
得到的是 $z-v^*(z,1)=\mathbb E[X_0]$，所有样本都坍缩为数据均值。
这与从整个时间区间积分得到正确端点分布并不矛盾。

Rectified Flow 的 reflow 再用已有流把噪声与生成端点配成新对，
在新配对上重新进行直线回归，希望减少轨迹弯曲与少步误差。
它改变的是端点耦合；独立配对、模型产生的配对与最优传输配对不是同一个东西。
二次代价最优传输寻找使 $\mathbb E\|Z-X_0\|^2$ 最小的合法耦合，
而一个小批量内的最优匹配只近似这个全分布问题。
Schrödinger Bridge 又保留了随机参考过程，在端点约束下最小化路径 KL；
它与确定性运输有联系，也不能仅凭“连接两端分布”就与 Rectified Flow 等同。
这些区别决定了训练配对、路径随机性与实际采样的含义。\cite{theorybook}
